
\section{Implementation Details}
\subsection{Model Details}
For the task of sampling from $x \in p_\theta(x|y_0)$ where $x \in \mathcal{X}, y_0 \in \mathcal{Y}^n$, the model will first sample from an isotropic Gaussian distribution $z_T \in \mathbb{R}^{C'\times H'\times W'} \sim p_\theta(z_{T}|y_0)$, and perform a $T$-step denoising procedure to gradually produce less noisy samples $z_{T-1} \sim p_\theta(z_{T-1}|y_0), \cdots, z_{0} \sim p_\theta(z_{0}|y_0)$. $C',H',W'$ refer to the latent dimension of the Stable Diffusion model. After reaching timestep $0$, the model will then obtain the final output image $x = \mathcal{D}_\theta(z_0)$ by mapping the resulting latent representation $z_0 \in \mathbb{R}^{C'\times H'\times W'}$ to the image space via a pretrained VQGAN~\citep{vqgan} quantization-decoder $\mathcal{D}_\theta$.

We use the official implementations and Stable Diffusion~\citep{stablediffusion} as the backbone model for all models. Based on the open source availability, we choose Stable Diffusion 2.1 as the base model for all methods except StructureDiffusion~\citep{structurediffusion} which uses Stable Diffusion 1.4, and GLIGEN~\citep{gligen} which adds additional modules to the model. For GLIGEN, we use the open sourced model pretrained on Flickr~\citep{flickr}, VG~\citep{vg}, Object365~\citep{object365}, SBU~\citep{sbu} and CC3M~\citep{cc3m} dataset. For datasets with incomplete localization information, they generate pseudo captions with class names and pseudo bounding boxes with GLIP~\citep{glip}, which is a open-vocabulary semantic rich object detector. 

We apply our cross attention control method to one model from each category in order to demonstrate the effectiveness of our method in diverse circumstances. We denote our method as CAC and integrate our method to Stable Diffusion 2.1, MultiDiffusion~\citep{multidiffusion} and GLIGEN as an add-on. Notice that before applying our technique, only MultiDiffusion and GLIGEN have the ability to process location information. All images are generated at $512\times 512$ resolution, which corresponds to $C = 3, H = 512, W = 512$. The latent space of VQGAN has dimensions $C' = 4, H' = 64, W' = 64$. Notice that $\lambda$ is a hyperparameter that one should tune in practice. For all quantitative analysis, we use $\lambda = 1$ for all the captions and $\lambda = 10$ for all the localized prompts. For experiments that incorporate CAC with MultiDiffusion, We use MD ratio of $0.4$ to report the quantitative results for both COCO and Cityscapes experiments and use $0.25$ for the CC-500 experiment. However, we do find lower MD ratios and $\lambda$ also provide satisfactory qualitative results. Details about each experiment setting are discussed in Section 4.

\subsection{Evaluation Details}
To evaluate the model performance, we use a variety of metrics for different aspects of the experiments. For fidelity, we use Kernel Inception Score (KID)~\citep{kid} between the generated images and filtered COCO validation set.
KID imposes fewer assumptions on the data distributions than FID~\citep{fid}, converges to the true values even with small numbers of samples and has also been widely used to evaluate fideilty of generative models~\citep{sdedit}.

For controllability, we use an YOLOv8 model~\cite{yolo} trained on the COCO dataset to predict the bounding boxes in the generated images and then calculate the precision (P), recall (R), mAP50
which is the mean average precision for bounding boxes that have an above 50\% intersection-over-union (IoU) with the ground truth boxes, 
and mAP50-95 
which is the mean average precision for bounding boxes that have a 50\% to 95\% IoU 
with the ground truth boxes.

Since the divergence between data distribution of Cityscapes images and the generated images is substantial according to the KID results in Table 2, semantic segmentation models that are only trained on Cityscapes will not work well in our setting due to the data distribution shift. As a result, we use Semantic Segment Anything (SSA)~\citep{ssa}, which is a general purpose open-vocabulary model that leverages Segment Anything Model (SAM)~\citep{sam} for semantic segmentation. We report the mean IoU (mIoU) score and mean accuracy (mACC) calculated from all classes, and all pixel accuracy (aACC).

In the compositional generation task we following~\citet{structurediffusion} and categorize each generation into three categories: generations that have incorrect or missing objects, generations that have the correct objects but the wrong colors, and the generations that have the correct objects with the correct colors. Similar to~\citet{gligen,structurediffusion}, we choose GLIP~\citep{glip}, the semantic rich open-vocabulary object detector to produce object and color detection results. We report the mean percentage of GLIP detection (mGLIP) with confidence threshold from 0.6 to 0.8 with 0.05 step size. Considering the difficulty in recognizing compositional objects, we also perform human evaluation on Amazon Mechanical Turk (mTurk) to verify the automatic results. Details about the human evaluation are in Appendix~\ref{app:human}.

We use the CC-500 dataset~\citep{structurediffusion} for the task of generating with compositional prompts. CC-500 consists of text prompts with the format ``a <color 1> <object 1> and a <color 2> <object 2>''. For methods that require location inputs, we apply a simple heuristic where the method will generate the first object on the left hand side of the image and the second object on the right hand side.
In other words, the first object will have a bounding box that spans the left half of the image and the second object will have one spanning the right half. All bounding boxes leave a 40-pixel margin to each border and the middle line of the image.

COCO is realised under Creative Commons Attribution 4.0 License. Cityscapes, CC-500 and GLIP are realised under MIT License. Yolov8 is realised under GNU Affero General Public License v3.0. SAM and SSA are realised under Apache License 2.0.  We use the default thresholds for all metrics.